

\documentclass[notes, 11pt]{beamer} % THIS
% \documentclass[notes, 11pt]{beamer}
% \documentclass[notes, 11pt, aspectratio=169]{beamer}

% \usepackage{helvet}
% \usepackage[default]{lato}
% \usepackage{array}


% \documentclass[handout]{beamer}
% \usepackage{pgfpages}
% \pgfpagesuselayout{4 on 1}[a4paper,border shrink=5mm]

\usepackage{bbm}
\usepackage{amsthm}
\usepackage{amssymb}
\usepackage{xcolor}
\usepackage{graphicx}
\usepackage{epstopdf}
\usepackage{booktabs}
\usepackage[percent]{overpic}

\usepackage{sgamevar, tikz} % Game theory packages
\usetikzlibrary{trees, calc} % For extensive form games
\usepackage{subfig} % Manipulation and reference of small or sub figures and tables

\setbeamertemplate{theorems}[numbered] % to number
% \usetheme{Malmoe}
\usetheme{default}
% \usecolortheme{whale}
\usecolortheme{rose}


\newtheorem{defi}{Definition}
\newtheorem{teo}{Theorem}
\newtheorem{remi}{Remark}
\newtheorem{exi}{Example}
\newtheorem{propi}{Proposition}
\newtheorem{proteo}{Proof}
\newtheorem{lemi}{Lemma}
\newtheorem{prolemi}{Proof}
\newtheorem{propropi}{Proof}
\newtheorem{coro}{Corollary}
\newtheorem{proci}{Algorithm}
\newtheorem{ass}{Assumption}

\setbeamertemplate{footline}[frame number]

\usepackage{caption}
\captionsetup[figure]{labelformat=empty}%

\newenvironment{wideitemize}{\itemize\addtolength{\itemsep}{10pt}}{\enditemize}
% \usepackage{enumitem}

\usepackage[spanish]{babel}

% \renewcommand{\baselinestretch}{1}
% \expandafter\def\expandafter\insertshorttitle\expandafter{%
%   \insertshorttitle\hfill%
%   \insertframenumber\,/\,\inserttotalframenumber}

\DeclareMathOperator*{\argmax}{arg\,max}

\begin{document}
\title{Microeconom\'ia Aplicada (MAE)}   
\subtitle{Trabajo de Investigaci\'on}   
\author[Paz y Mino ]{Jos\'e Manuel Paz y Mi\~no}
\institute[shortinst]{FEN UCHILE}
\date{Oto\~no 2025} 

\frame{\titlepage} 

\begin{frame}{Objetivo General}
	\begin{wideitemize}
		\item<1-> Manejar modelos te\'oricos de \'area de interes.
		\item<2-> Idealmente, modelos te\'oricos usando en el trabajo de investigaci\'on sirven para tesis de magister y/o AFE.
	\end{wideitemize}
\end{frame}

\begin{frame}{Requisitos}
	\begin{wideitemize}
		\item<1-> De 1 o 2 personas.
		\item<2-> Es un trabajo te\'orico, no emp\'irico (no datos).
		\item<3-> Tiene 2 entregas
		\begin{wideitemize}
			\item<4->[--] Primera entrega (Mayo 14): exposici\'on de modelo te\'orico ya publicado relacionado a \'area de interes.
			\item<5->[--] Segunda entrega (Julio 1): extensi\'on de modelo te\'orico ya publicado o planteamiento de un modelo nuevo.
		\end{wideitemize}
	\end{wideitemize}
\end{frame}

\begin{frame}{Sobre primera entrega}
	\begin{wideitemize}
		\item<1-> Elecci\'on de area de interes.
		\item<2-> B\'usqueda de investigaci\'on relevante.
		\item<3-> Leer investigaci\'on relevante.
		\item<4-> Exponer investigaci\'on relevante.
	\end{wideitemize}
\end{frame}

\begin{frame}{Sobre primera entrega}{Elecci\'on de area de interes}
	\begin{wideitemize}
		\item<1-> Referido a una tema: pol\'itica, salud, laboral, etc.
		\item<2-> Basado en preferencias individuales.
		\item<3-> Puede ser un buen momento para conversar con profesores de la FEN.
		\item<4-> En mi caso, mi \'area de interes es organizaci\'on industrial y pol\'iticas de competencia.
	\end{wideitemize}
\end{frame}

\begin{frame}{Sobre primera entrega}{B\'usqueda de investigaci\'on relevante}
	\begin{wideitemize}
		\item<1-> En cada area de interes hay innumerables sub-areas.
		\begin{wideitemize}
			\item<2->[--] Revisar reviews: Annual Review of Economics, Journal of Economic Literature, Journal of Economic Perspectives, Annual Review of Resource Economics, Handbook of XYZ Economics.
			\item<3->[--] Revisar libros: Biblioteca FEN.
			\item<4->[--] Si uno tiene un tema espec\'ifico revisar directamente en Google Scholar.
		\end{wideitemize}
		\item<5-> Muchas veces investigaciones recientes usan m\'etodos avazandos. Pero preguntas importantes tienen larga historia. Buscar referencias m\'as antiguas si son complicadas de leer.
	\end{wideitemize}
\end{frame}

\begin{frame}{Sobre primera entrega}{Leer investigaci\'on relevante}
	\begin{wideitemize}
		\item<1-> Identificar las secci\'on de la investigaci\'on.
		\begin{wideitemize}
			\item<2->[--] Introducci\'on/Motivaci\'on, Modelo, Est\'atica Comparativa, Extensiones, Conclusiones.
		\end{wideitemize}
		\item<3-> Para poder entender la investigacion leer varias veces la introducci\'on/motivacion y conclusiones.
		\item<4-> Una buena investigaci\'on es clara en su contribuci\'on.
	\end{wideitemize}
\end{frame}

\begin{frame}{Sobre primera entrega}{Leer investigaci\'on relevante}
	\begin{wideitemize}
		\item<1-> Leer el modelo.
		\begin{wideitemize}
			\item<2->[--] Elementos del modelo.
			\item<3->[--] Lemmas, proposiciones.
			\item<4->[--] Revisar demostraciones.
		\end{wideitemize}
	\end{wideitemize}
\end{frame}

\begin{frame}{Sobre primera entrega}{Exponer investigaci\'on relevante}
	\begin{wideitemize}
		\item<1-> Se busca que el alumno/a
		\begin{wideitemize}
			\item<2->[--] Presente el modelo.
			\item<3->[--] Presente lemmas, proposiciones relevantes (discusi\'on intuitiva).
			\item<4->[--] Est\'atica comparativa.
		\end{wideitemize}
	\end{wideitemize}
\end{frame}


\begin{frame}
\maketitle
\end{frame}

\end{document}
